\section{RQ4: What Do Deterministic Reductions Save?}
\label{sec:rq4}

Table~\ref{tab:arms} lists the arms over the \NxReplayed{} replayed networkx edits.
Benchmarks are excluded from these totals; RQ1 shows that they would dominate them.

\begin{table}[t]
\caption{Arms over the replayed networkx edits: tests run, estimated and measured end-to-end test cost (build plus tests plus start-up, no benchmarks), and misses of the frozen functional mutants (Wilson 95\pct{} CI of the miss rate). Estimated seconds come from per-test durations and are comparable across arms; measured seconds are single runs under different load and are not comparable with the estimates. INC equals FULL in tests run; STATIC-BENCH equals TESTMON in tests and differs only in benchmarks.}
\label{tab:arms}
\centering\footnotesize\setlength{\tabcolsep}{4pt}
\begin{tabular}{@{}lrrrl@{}}
\toprule
Arm & Tests & Est.\ s & Meas.\ s & Missed (CI) \\
\midrule
FULL (A0) & \TestsRunFull & \EstFull & \FullMeasured & \MutantsNone/\MutantsKilled{} \ci{\MutantsNone}{\ZeroMissHi} \\
TESTMON (A2) & \TestsRunTestmon & \EstTestmon & \TestmonMeasured & \MutantsNone/\MutantsKilled{} \ci{\MutantsNone}{\ZeroMissHi} \\
CHANGED-MOD.\ (A3) & \TestsRunAThree & \EstAThree & -- & \AThreeMissed/\MutantsKilled{} \ci{\AThreeMissLo}{\AThreeMissHi} \\
A4 (ours) & \TestsRunAFour & \EstAFour & -- & \MutantsNone/\MutantsKilled{} \ci{\MutantsNone}{\ZeroMissHi} \\
\bottomrule
\end{tabular}
\end{table}

\paragraph*{Incremental build}
For networkx, an incremental build saves almost nothing: the clean build is \NxBuildShareTestsLo--\NxBuildShareTestsHi\pct{} of the test time.
For bottleneck the build is a larger share of build plus tests (Section~\ref{sec:rq1}), but small once benchmarks count.
We report it as negligible for this project, not as a saving.

\paragraph*{Test selection, end to end}
pytest-testmon, unchanged, took \TestmonMeasured\,s over the replayed sequence against \FullMeasured\,s for full serial tests and builds, and \FullIdle\,s in the idle re-run: \TestmonRatioLo--\TestmonRatioHi\x{} cheaper.
Counting only the seconds of the tests run would suggest \TestSecRatio\x{} (\TestmonTestSec{} against \FullTestSec\,s).
The difference is pytest start-up, paid per step, and testmon's own analysis; its first run, which builds the database, took \TestmonDB\,s against \TestmonDBPlain\,s for a plain run.
The testmon measurement ran concurrently with other work and the full run includes one noisy state, so the ratio is a range.
On the worked example of step~10, test time falls from about \WorkedTestsBefore\,s to \WorkedTestsAfter\,s, about \WorkedRatio\x, because start-up remains.
Reductions in tests run, as optimizers report them (PerfAgent reports large shares of tests skipped~\cite{deng2026perfagent}), are therefore not proportional reductions in validation time: on networkx, a \TestSecRatio\x{} cut in test-seconds became \TestmonRatioLo--\TestmonRatioHi\x{} end to end, in one run under mixed load.

\paragraph*{Functional safety}
Every arm except CHANGED-MODULE caught all \MutantsKilled{} frozen mutants at the step that introduced them.
CHANGED-MODULE missed \AThreeMissed: in both, the killing tests live in a test file whose name does not match the module.
With a full check every three or every five edits, both are caught one or two steps later.
Because the mutants sit in the edited functions, every arm that selects by dependency catches them by construction; this fault model cannot rank selection rules, only show that the cheap arm is not obviously unsafe.

\paragraph*{Our accumulated-change policy (negative result)}
A4 selects tests over the change accumulated since its last full check and triggers a full check when an additive change score crosses a fixed threshold.
On this sequence it behaved like periodic full checks every \AFourK{} edits or so, and it bought no extra detection: plain testmon catches the same mutants at about half its estimated cost (\EstTestmon{} against \EstAFour\,s).
The accumulation null, the collapse of its benchmark selection and the batching comparison are in Appendix~\ref{app:a4}.

\paragraph*{Answer to RQ4}
On networkx, existing test selection saves \TestmonRatioLo--\TestmonRatioHi\x{} of test-validation time end to end with no functional miss on the frozen mutants, and incremental builds save almost nothing.
Our deterministic accumulated-change policy adds nothing measurable under this fault model.
Neither touches the benchmark cost of RQ1.
